\documentclass[UTF8,11pt]{ctexart}
\usepackage[a4paper,margin=25mm]{geometry}
\usepackage{amsmath,amssymb,amsthm,mathtools}
\usepackage[all]{xy}
\usepackage{xurl}
\usepackage{hyperref}
\usepackage{enumitem}
\hypersetup{hidelinks,breaklinks=true}
\setlength{\emergencystretch}{3em}
\newtheorem*{theorem}{Theorem}
\newtheorem*{lemma}{Lemma}
\newtheorem*{proposition}{Proposition}
\newtheorem*{corollary}{Corollary}
\newtheorem*{definition}{Definition}
\newcommand{\Spec}{\operatorname{Spec}}
\newcommand{\Hom}{\operatorname{Hom}}
\newcommand{\End}{\operatorname{End}}
\newcommand{\Gal}{\operatorname{Gal}}
\newcommand{\Cl}{\operatorname{Cl}}
\newcommand{\vol}{\operatorname{vol}}
\newcommand{\covol}{\operatorname{covol}}
\newcommand{\tr}{\operatorname{tr}}
\newcommand{\im}{\operatorname{im}}
\newcommand{\id}{\operatorname{id}}
\newcommand{\coker}{\operatorname{coker}}
\newcommand{\stacksref}[2]{\href{https://stacks.math.columbia.edu/tag/#1}{#2 [#1]}}
\begin{document}
\section*{004\quad Noetherian环的Serre正规性判据}
\subsection*{来源与整理范围}
The Stacks Project Authors, \emph{The Stacks Project}, Chapter 10, \emph{Commutative Algebra}。主命题：Lemma 10.157.4，Serre's criterion for normality，Tag 031S。收录原定义 10.157.1 中的环条件，以及所用 Lemma 10.157.3（031R）和 Lemma 10.119.2（0BHZ）的作者证明。后者原页面归因于 J\'anos Koll\'ar, \emph{Variants of normality for Noetherian schemes}。
\par\url{https://stacks.math.columbia.edu/tag/031S}
\par\url{https://stacks.math.columbia.edu/tag/031O}
\par\url{https://stacks.math.columbia.edu/tag/031R}
\par\url{https://stacks.math.columbia.edu/tag/0BHZ}
\par 使用版本：2026-09-28 实际访问的在线正文；未取得该网页构建的固定 commit，不编造版本号。原文数学公式在网页中以 TeX 提供；本条保留其内容，引用按所访问页面显示的编号列出。

\subsection*{作者原命题与人类证明}
\begin{definition}[10.157.1, ring conditions]
Let $R$ be a Noetherian ring. Let $k\geq0$ be an integer.
\begin{enumerate}
\item We say $R$ has property $(R_k)$ if for every prime $\mathfrak p$ of height $\leq k$ the local ring $R_{\mathfrak p}$ is regular. We also say that $R$ is regular in codimension $\leq k$.
\item We say $R$ has property $(S_k)$ if for every prime $\mathfrak p$ the local ring $R_{\mathfrak p}$ has depth at least $\min\{k,\dim(R_{\mathfrak p})\}$.
\end{enumerate}
\end{definition}

\begin{lemma}[10.119.2 (Koll\'ar), Tag 0BHZ]
Let $(R,\mathfrak m)$ be a local Noetherian ring. Then exactly one of the following holds:
\begin{enumerate}
\item $(R,\mathfrak m)$ is Artinian,
\item $(R,\mathfrak m)$ is regular of dimension $1$,
\item $\operatorname{depth}(R)\geq2$, or
\item there exists a finite ring map $R\to R'$ which is not an isomorphism whose kernel and cokernel are annihilated by a power of $\mathfrak m$ such that $\mathfrak m$ is not an associated prime of $R'$ and $R'\ne0$.
\end{enumerate}
\end{lemma}
\begin{proof}
Observe that $(R,\mathfrak m)$ is not Artinian if and only if $V(\mathfrak m)\subset\Spec(R)$ is nowhere dense. See Proposition 10.60.7. We assume this from now on.

Let $J\subset R$ be the largest ideal killed by a power of $\mathfrak m$. If $J\ne0$ then $R\to R/J$ shows that $(R,\mathfrak m)$ is as in (4).

Otherwise $J=0$. In particular $\mathfrak m$ is not an associated prime of $R$ and we see that there is a nonzerodivisor $x\in\mathfrak m$ by Lemma 10.63.18. If $\mathfrak m$ is not an associated prime of $R/xR$ then $\operatorname{depth}(R)\geq2$ by the same lemma. Thus we are left with the case when there is a $y\in R$, $y\notin xR$ such that $y\mathfrak m\subset xR$.

If $y\mathfrak m\subset x\mathfrak m$ then we can consider the map $\varphi:\mathfrak m\to\mathfrak m$, $f\mapsto yf/x$ (well defined as $x$ is a nonzerodivisor). By the determinantal trick of Lemma 10.16.2 there exists a monic polynomial $P$ with coefficients in $R$ such that $P(\varphi)=0$. We conclude that $P(y/x)=0$ in $R_x$. Let $R'\subset R_x$ be the ring generated by $R$ and $y/x$. Then $R\subset R'$ and $R'/R$ is a finite $R$-module annihilated by a power of $\mathfrak m$. Thus $R$ is as in (4).

Otherwise there is a $t\in\mathfrak m$ such that $yt=ux$ for some unit $u$ of $R$. After replacing $t$ by $u^{-1}t$ we get $yt=x$. In particular $y$ is a nonzerodivisor. For any $t'\in\mathfrak m$ we have $yt'=xs$ for some $s\in R$. Thus $y(t'-st)=xs-xs=0$. Since $y$ is not a zero-divisor this implies that $t'=ts$ and so $\mathfrak m=(t)$. Thus $(R,\mathfrak m)$ is regular of dimension $1$.

The argument so far shows that every $R$ falls into one of the $4$ cases. To finish we have to show that no two of the properties can hold simultaneously for each of the $6$ combinations of properties. We'll just show that (2) and (3) each exclude (4); the other cases are left to the reader.

Assume $R$ is regular of dimension $1$ and that $R\to R'$ is a finite ring map whose kernel and cokernel are annihilated by a power of $\mathfrak m$ and $\mathfrak m$ is not an associated prime of $R'$. Then $R'$ is a finite $R$-module of depth at least $1$ and hence free by Lemma 10.106.6. Since $R\to R'$ is an isomorphism in the generic point, we conclude that $R'$ must be free of rank $1$ as an $R$-module. Then $R'\to\End_R(R')\cong R$ is an inverse to the map $R\to R'$. Thus (4) cannot be true.

Assume $R$ has depth $\geq2$ and that $R\to R'$ is a finite ring map whose kernel and cokernel are annihilated by a power of $\mathfrak m$ and $\mathfrak m$ is not an associated prime of $R'$. Then $R\to R'$ is necessarily injective (as the kernel would have both depth $0$ and depth $\geq1$) and the cokernel has depth at least $1$ (by Lemma 10.72.6 and the fact that $R'$ has depth $\geq1$ by assumption) whence cannot be supported on $\{\mathfrak m\}$ unless it is $0$ as well. Thus (4) cannot be true.
\end{proof}

\begin{lemma}[10.157.3, Tag 031R]
Let $R$ be a Noetherian ring. The following are equivalent:
\begin{enumerate}
\item $R$ is reduced, and
\item $R$ has properties $(R_0)$ and $(S_1)$.
\end{enumerate}
\end{lemma}
\begin{proof}
Suppose that $R$ is reduced. Then $R_{\mathfrak p}$ is a field for every minimal prime $\mathfrak p$ of $R$, according to Lemma 10.25.1. Hence we have $(R_0)$. Let $\mathfrak p$ be a prime of height $\geq1$. Then $A=R_{\mathfrak p}$ is a reduced local ring of dimension $\geq1$. Hence its maximal ideal $\mathfrak m$ is not an associated prime since this would mean there exists an $x\in\mathfrak m$ with annihilator $\mathfrak m$ so $x^2=0$. Hence the depth of $A=R_{\mathfrak p}$ is at least one, by Lemma 10.63.9. This shows that $(S_1)$ holds.

Conversely, assume that $R$ satisfies $(R_0)$ and $(S_1)$. If $\mathfrak p$ is a minimal prime of $R$, then $R_{\mathfrak p}$ is a field by $(R_0)$, and hence is reduced. If $\mathfrak p$ is not minimal, then we see that $R_{\mathfrak p}$ has depth $\geq1$ by $(S_1)$ and we conclude there exists an element $t\in\mathfrak pR_{\mathfrak p}$ such that $R_{\mathfrak p}\to R_{\mathfrak p}[1/t]$ is injective. Now $R_{\mathfrak p}[1/t]$ is contained in the product of its localizations at prime ideals, see Lemma 10.23.1. This implies that $R_{\mathfrak p}$ is a subring of a product of localizations of $R$ at $\mathfrak p\supset\mathfrak q$ with $t\notin\mathfrak q$. Since these primes have smaller height by induction on the height we conclude that $R$ is reduced.
\end{proof}

\begin{lemma}[10.157.4 (Serre's criterion for normality), Tag 031S]
Let $R$ be a Noetherian ring. The following are equivalent:
\begin{enumerate}
\item $R$ is a normal ring, and
\item $R$ has properties $(R_1)$ and $(S_2)$.
\end{enumerate}
\end{lemma}
\begin{proof}
Proof of (1) $\Rightarrow$ (2). Assume $R$ is normal, i.e., all localizations $R_{\mathfrak p}$ at primes are normal domains. In particular we see that $R$ has $(R_0)$ and $(S_1)$ by Lemma 10.157.3. Hence it suffices to show that a local Noetherian normal domain $R$ of dimension $d$ has depth $\geq\min(2,d)$ and is regular if $d=1$. The assertion if $d=1$ follows from Lemma 10.119.7.

Let $R$ be a local Noetherian normal domain with maximal ideal $\mathfrak m$ and dimension $d\geq2$. Apply Lemma 10.119.2 to $R$. It is clear that $R$ does not fall into cases (1) or (2) of the lemma. Let $R\to R'$ as in (4) of the lemma. Since $R$ is a domain we have $R\subset R'$. Since $\mathfrak m$ is not an associated prime of $R'$ there exists an $x\in\mathfrak m$ which is a nonzerodivisor on $R'$. Then $R_x=R'_x$ so $R$ and $R'$ are domains with the same fraction field. But finiteness of $R\subset R'$ implies every element of $R'$ is integral over $R$ (Lemma 10.36.3) and we conclude that $R=R'$ as $R$ is normal. This means (4) does not happen. Thus we get the remaining possibility (3), i.e., $\operatorname{depth}(R)\geq2$ as desired.

Proof of (2) $\Rightarrow$ (1). Assume $R$ satisfies $(R_1)$ and $(S_2)$. By Lemma 10.157.3 we conclude that $R$ is reduced. Hence it suffices to show that if $R$ is a reduced local Noetherian ring of dimension $d$ satisfying $(S_2)$ and $(R_1)$ then $R$ is a normal domain. If $d=0$, the result is clear. If $d=1$, then the result follows from Lemma 10.119.7.

Let $R$ be a reduced local Noetherian ring with maximal ideal $\mathfrak m$ and dimension $d\geq2$ which satisfies $(R_1)$ and $(S_2)$. By Lemma 10.37.16 it suffices to show that $R$ is integrally closed in its total ring of fractions $Q(R)$. Pick $x\in Q(R)$ which is integral over $R$. Then $R'=R[x]$ is a finite ring extension of $R$ (Lemma 10.36.5). Because $\dim(R_{\mathfrak p})<d$ for every nonmaximal prime $\mathfrak p\subset R$ we have $R_{\mathfrak p}=R'_{\mathfrak p}$ by induction. Hence the support of $R'/R$ is $\{\mathfrak m\}$. It follows that $R'/R$ is annihilated by a power of $\mathfrak m$ (Lemma 10.62.4). By Lemma 10.119.2 this contradicts the assumption that the depth of $R$ is $\geq2=\min(2,d)$ and the proof is complete.
\end{proof}


\subsection*{整理说明（非作者正文）}
本条只计 Serre 正规性判据一个问题。0BHZ 与 031R 是其证明的支撑，不另计数量。0BHZ 原文明确把部分互斥情形留给读者，本转录保留该说明；用于本主命题的四分情形构造以及 (3) 排除 (4) 的关键论证，作者已写出且均在本文中，未以“留作练习”跳过主命题的核心步骤。

标准先修保留原引用：一维 Noetherian 局部正规整环与离散赋值环的对应、有限整扩张基本性质、伴随素理想与非零因子、深度引理、整元素的行列式技巧和总分式环的基本性质。主命题的两个方向均完整收录。网页评论不是证明正文，未混入作者论证。按原编号保留引用，而不是使用未定义的内部 \verb|\ref|。原生数学 TeX 转成独立文档；没有取得整个原始章节文件，保存的是所用正文摘录。

\subsection*{可修改的简短标注（非作者正文）}
$c$：Noetherian 环的正规性判别。

$a$：素理想局部化；余维正则性；深度与非零因子；伴随素理想；有限整扩张；总分式环；行列式技巧。

\texttt{cross\_theory\_status = uncertain}。当前可见的联系是把整闭性转成局部深度和余维条件，关键位置为 0BHZ 的有限扩张构造与 031S 的维数归纳；本轮不进一步认证其跨理论程度。全部标注为非穷尽、可修改初稿。
\subsection*{许可与修改记录}
原数学正文作者：The Stacks Project Authors（集体作者）。原文按 GNU Free Documentation License 1.2 或后续版本发布；本摘录的整理部分随原文采用相同许可。许可全文随包保存在 \texttt{sources/stacks/GFDL-1.2.txt}。无不变章节、封面文字或封底文字。本文题名与来源作品题名不同，不暗示原作者认可本次选题或标注。

\textbf{History.} 2026-09-28：ChatGPT 为本对话材料包整理摘录，加入题号、来源定位、独立排版与简短标注；数学正文的具体排版变动见整理说明。

\end{document}
